\Needspace{8\baselineskip}
\Article{188}
Da alle Formen derselben Periode eigentlich äquivalent sind, stellt sich die Frage, ob auch Formen aus verschiedenen Perioden eigentlich äquivalent sein können. Bevor wir aber zeigen, dass \textit{dies unmöglich ist}, müssen wir einige Punkte über die Transformation reduzierter Formen behandeln.

Da wir im Folgenden sehr oft mit Transformationen von Formen zu tun haben, verwenden wir von jetzt an stets die folgende abgekürzte Schreibweise, um die Schreibweise möglichst kurz zu halten. Wird eine Form $LXX+2MXY+NYY$ durch die Substitution $X=\alpha x+\beta y$, $Y=\gamma x+\delta y$ in die Form $lxx+2mxy+nyy$ transformiert, so sagen wir einfach, $(L,M,N)$ werde in $(l,m,n)$ durch die Substitution $\alpha,\beta,\gamma,\delta$ transformiert. So brauchen wir die Unbestimmten der einzelnen betrachteten Formen nicht mit eigenen Zeichen zu bezeichnen. --- Dabei muss aber die \textit{erste} Unbestimmte in jeder Form sorgfältig von der \textit{zweiten} unterschieden werden.

Gegeben sei die reduzierte Form $(a,b,-a')\mathinner{\ldotp\ldotp\ldotp} f$ mit Determinante $D$. Wie in Artikel~186 bilde man die nach beiden Seiten unendliche Folge reduzierter Formen $\mathinner{\ldotp\ldotp\ldotp\ldotp} {}''f,{}'f,f,f',f''\mathinner{\ldotp\ldotp\ldotp\ldotp}$, und zwar sei
\SourceLeaf{177}
\[
\begin{aligned}
f'&=(-a',b',a''),& f''&=(a'',b'',-a''')\quad\text{usw.}\\
{}'f&=(-{}'a,{}'b,a),& {}''f&=({}''a,{}''b,-{}'a)\quad\text{usw.}
\end{aligned}
\]
Man setze
\[
\begin{array}{llll}
\dfrac{b+b'}{-a'}=h',&\dfrac{b'+b''}{a''}=h'',&\dfrac{b''+b'''}{-a'''}=h'''&\text{usw.}\\[6pt]
\dfrac{{}'b+b}{a}=h,&\dfrac{{}''b+{}'b}{-{}'a}={}'h,&\dfrac{{}'''b+{}''b}{{}''a}={}''h&\text{usw.}
\end{array}
\]
Dann wird offenbar, wenn die Zahlen $\alpha',\alpha'',\alpha'''$ usw., $\beta',\beta'',\beta'''$ usw. und so weiter (wie in Artikel~177) nach dem folgenden Algorithmus gebildet werden,
{\small
\setlength{\arraycolsep}{3pt}
\[
\begin{array}{llll}
\alpha'=0&\beta'=-1&\gamma'=1&\delta'=h'\\
\alpha''=\beta'&\beta''=h''\beta'&\gamma''=\delta&\delta''=h''\delta'-1\\
\alpha'''=\beta''&\beta'''=h'''\beta''-\beta'&\gamma'''=\delta''&\delta'''=h'''\delta''-\delta'\\
\alpha''''=\beta'''&\beta''''=h''''\beta'''-\beta''&\gamma''''=\delta'''&\delta''''=h''''\delta'''-\delta''
\\
\multicolumn{4}{c}{\text{usw.}}
\end{array}
\]
}
$f$ transformiert
\[
\begin{array}{@{}rlll@{}}
\text{in}&f'&\text{durch die Substitution}&\alpha',\ \beta',\ \gamma',\ \delta'\\
&f''&\mathinner{\ldotp\ldotp\ldotp\ldotp\ldotp\ldotp\ldotp\ldotp}&\alpha'',\ \beta'',\ \gamma'',\ \delta''\\
&f'''&\mathinner{\ldotp\ldotp\ldotp\ldotp\ldotp\ldotp\ldotp\ldotp}&\alpha''',\ \beta''',\ \gamma''',\ \delta'''
\end{array}
\]
\centerline{usw.}
und alle diese Transformationen sind eigentlich.

\Needspace{24\baselineskip}
Da ${}'f$ in $f$ durch die eigentliche Substitution $0,-1,1,h$ übergeht (Artikel~158), geht $f$ in ${}'f$ durch die eigentliche Substitution $h,1,-1,0$ über. Aus ähnlichem Grund geht ${}'f$ in ${}''f$ durch die eigentliche Substitution ${}'h,1,-1,0$ über, ${}''f$ in ${}'''f$ durch die eigentliche Substitution ${}''h,1,-1,0$ usw. Daraus folgt nach Artikel~159 auf dieselbe Weise wie in Artikel~177: Werden die Zahlen ${}'\alpha,{}''\alpha,{}'''\alpha$ usw., ${}'\beta,{}''\beta,{}'''\beta$ usw. und so weiter nach dem folgenden Algorithmus gebildet,
{\small
\setlength{\arraycolsep}{3pt}
\[
\begin{array}{llll}
{}'\alpha=h&{}'\beta=1&{}'\gamma=-1&{}'\delta=0\\
{}''\alpha={}'h\,{}'\alpha-1&{}''\beta={}'\alpha&{}''\gamma={}'h\,{}'\gamma&{}''\delta={}'\gamma\\
{}'''\alpha={}''h\,{}''\alpha-{}'\alpha&{}'''\beta={}''\alpha&{}'''\gamma={}''h\,{}''\gamma-{}'\gamma&{}'''\delta={}''\gamma\\
{}''''\alpha={}'''h\,{}'''\alpha-{}''\alpha&{}''''\beta={}'''\alpha&{}''''\gamma={}'''h\,{}'''\gamma-{}''\gamma&{}''''\delta={}'''\gamma
\\
\multicolumn{4}{c}{\text{usw.}}
\end{array}
\]
}
so wird $f$ transformiert
\SourceLeaf{178}
\[
\begin{array}{@{}rlll@{}}
\text{in}&{}'f&\text{durch die Substitution}&{}'\alpha,\ {}'\beta,\ {}'\gamma,\ {}'\delta\\
&{}''f&\mathinner{\ldotp\ldotp\ldotp\ldotp\ldotp\ldotp\ldotp\ldotp}&{}''\alpha,\ {}''\beta,\ {}''\gamma,\ {}''\delta\\
&{}'''f&\mathinner{\ldotp\ldotp\ldotp\ldotp\ldotp\ldotp\ldotp\ldotp}&{}'''\alpha,\ {}'''\beta,\ {}'''\gamma,\ {}'''\delta
\end{array}
\]
\centerline{usw.}
und alle diese Transformationen sind eigentlich.

Setzt man $\alpha=1$, $\beta=0$, $\gamma=0$, $\delta=1$, so stehen diese Zahlen zur Form $f$ in derselben Beziehung wie $\alpha',\beta',\gamma',\delta'$ zu $f'$, $\alpha'',\beta'',\gamma'',\delta''$ zu $f''$ usw., ${}'\alpha,{}'\beta,{}'\gamma,{}'\delta$ zu ${}'f$ usw. Durch die Substitution $\alpha,\beta,\gamma,\delta$ geht nämlich die Form $f$ in $f$ über. Die unendlichen Folgen $\alpha',\alpha'',\alpha'''$ usw., ${}'\alpha,{}''\alpha,{}'''\alpha$ usw. lassen sich dann durch Einschieben des Gliedes $\alpha$ so verbinden, dass man sie als eine einzige zusammenhängende, nach beiden Seiten unendliche Folge auffassen kann, die überall demselben Gesetz folgt: $\mathinner{\ldotp\ldotp\ldotp} {}'''\alpha,{}''\alpha,{}'\alpha,\alpha,\alpha',\alpha'',\alpha'''\mathinner{\ldotp\ldotp\ldotp}$. Dieses Gesetz lautet
\[
\begin{gathered}
{}'''\alpha+{}'\alpha={}''h\,{}''\alpha,\quad {}''\alpha+\alpha={}'h\,{}'\alpha,\quad {}'\alpha+\alpha'=h\alpha,\\
\alpha+\alpha''=h'\alpha',\quad \alpha'+\alpha'''=h''\alpha''\text{ usw.}
\end{gathered}
\]
oder allgemein (wenn ein rechts geschriebener negativer Index dasselbe bedeuten soll wie ein links geschriebener positiver)
\[
\alpha^{m-1}+\alpha^{m+1}=h^m{\alpha^m}
\]
Ebenso ist die Folge $\mathinner{\ldotp\ldotp\ldotp} {}''\beta,{}'\beta,\beta,\beta',\beta''\mathinner{\ldotp\ldotp\ldotp}$ zusammenhängend; ihr Gesetz lautet
\[
\beta^{m-1}+\beta^{m+1}=h^{m+1}\beta^m
\]
und sie ist im Wesentlichen mit der vorigen identisch, wobei alle Glieder um eine Stelle verschoben sind: ${}''\beta={}'\alpha$, ${}'\beta=\alpha$, $\beta=\alpha'$ usw. Das Gesetz der zusammenhängenden Folge $\mathinner{\ldotp\ldotp\ldotp\ldotp} {}''\gamma,\allowbreak {}'\gamma,\allowbreak \gamma,\allowbreak \gamma',\allowbreak \gamma''\mathinner{\ldotp\ldotp\ldotp\ldotp\ldotp}$ lautet
\[
\gamma^{m-1}+\gamma^{m+1}=h^m{\gamma^m}
\]
und das Gesetz von $\mathinner{\ldotp\ldotp\ldotp} {}''\delta,{}'\delta,\delta,\delta',\delta''\mathinner{\ldotp\ldotp\ldotp}$ lautet
\[
\delta^{m-1}+\delta^{m+1}=h^{m+1}\delta^m
\]
außerdem gilt allgemein $\delta^m=\gamma^{m+1}$.

\textit{Beispiel.} Die gegebene Form $f$ sei $(3,8,-5)$. Sie wird transformiert
\SourceLeaf{179}
\[
\setlength{\arraycolsep}{4pt}
\begin{array}{@{}lll@{}}
\multicolumn{2}{l}{\text{in die Form}}&\text{durch die Substitution}\\
{}'''''''f&(-10,7,3)&-805,\ -152,\ +143,\ +27\\
{}''''''f&(3,8,-5)&-152,\ +45,\ +27,\ -8\\
{}'''''f&(-5,7,6)&+45,\ +17,\ -8,\ -3\\
{}''''f&(6,5,-9)&+17,\ -11,\ -3,\ +2\\
{}'''f&(-9,4,7)&-11,\ -6,\ +2,\ +1\\
{}''f&(7,3,-10)&-6,\ +5,\ +1,\ -1\\
{}'f&(-10,7,3)&+5,\ +1,\ -1,\ 0\\
f&(3,8,-5)&+1,\ 0,\ 0,\ +1\\
f'&(-5,7,6)&0,\ -1,\ +1,\ -3\\
f''&(6,5,-9)&-1,\ -2,\ -3,\ -7\\
f'''&(-9,4,7)&-2,\ +3,\ -7,\ +10\\
f''''&(7,3,-10)&+3,\ +5,\ +10,\ +17\\
f'''''&(-10,7,3)&+5,\ -8,\ +17,\ -27\\
f''''''&(3,8,-5)&-8,\ -45,\ -27,\ -152\\
f'''''''&(-5,7,6)&-45,\ +143,\ -152,\ +483
\end{array}
\]
\centerline{usw.}

\Needspace{8\baselineskip}
\Article{189}

Zu diesem Algorithmus ist Folgendes anzumerken.

1) Alle $a,a',a''$ usw., ${}'a,{}''a$ usw. haben dasselbe Vorzeichen; alle $b,b',b''$ usw., ${}'b,{}''b$ usw. sind positiv. In der Folge $\mathinner{\ldotp\ldotp\ldotp} {}''h,{}'h,h,h',h''\mathinner{\ldotp\ldotp\ldotp\ldotp}$ wechseln die Vorzeichen ab: Sind alle $a,a'$ usw. positiv, so ist $h^m$ oder ${}^m h$ für gerades $m$ positiv und für ungerades $m$ negativ; sind dagegen $a,a'$ usw. negativ, so ist $h^m$ oder ${}^m h$ für gerades $m$ negativ und für ungerades positiv.

2) Ist $a$ positiv und damit $h'$ negativ, $h''$ positiv usw., so ist $\alpha''=-1$ negativ, $\alpha'''=h''\alpha''$ negativ und $>\alpha''$ (oder $=\alpha''$, wenn $h''=1$); $\alpha''''=h'''\alpha'''-\alpha''$ positiv und $>\alpha'''$ (weil $h'''\alpha'''$ positiv und $\alpha''$ negativ ist); $\alpha'''''=h''''\alpha''''-\alpha'''$ positiv und $>\alpha''''$ (weil $h''''\alpha''''$ positiv ist) usw. Daraus folgt leicht, dass die Folge $\alpha',\alpha'',\alpha'''$ usw. ins Unendliche wächst und stets zwei positive Vorzeichen auf zwei negative folgen, sodass $\alpha^m$ das Vorzeichen $+,+,-,-$ hat, je nachdem $m\equiv0,1,2,3\ (\mathrm{mod.}\ 4)$. Ist $a$ negativ, so erhält man durch einen ähnlichen Schluss $\alpha''$ negativ, $\alpha'''$ positiv und entweder $>$ oder $=\alpha''$; $\alpha''''$ positiv und $>\alpha'''$; $\alpha'''''$ negativ und $>\alpha''''$ usw.,
\SourceLeaf{180}
sodass die Folge $\alpha',\alpha'',\alpha'''$ usw. fortwährend wächst und das Glied $\alpha^m$ das Vorzeichen $+,-,-,+$ hat, je nachdem $m\equiv0,1,2,3\ (\mathrm{mod.}\ 4)$.

\Needspace{32\baselineskip}
3) Auf diese Weise findet man, dass alle vier unendlichen Folgen $\alpha',\alpha'',\alpha'''$ usw., $\gamma,\gamma',\gamma''$ usw.; $\alpha',\alpha,{}'\alpha,{}''\alpha$ usw.; $\gamma,{}'\gamma,{}''\gamma$ usw. fortwährend wachsen, also auch die folgenden, mit ihnen identischen Folgen: $\beta,\beta',\beta''$ usw.; ${}'\delta,\delta,\delta',\delta''$ usw.; $\beta,{}'\beta,{}''\beta$ usw.; ${}'\delta,{}''\delta$ usw. Je nachdem $m\equiv0,1,2,3\ (\mathrm{mod.}\ 4)$, lauten die Vorzeichen
\[
\begin{array}{ll}
\text{von }\alpha^m,\quad +\ \pm\ -\ \mp;&\text{von }\beta^m,\quad \pm\ -\ \mp\ +\\[6pt]
\text{von }\gamma^m,\quad \pm\ +\ \mp\ -;&\text{von }\delta^m,\quad +\ \mp\ -\ \pm\\[6pt]
\text{von }{}^m\alpha,\quad +\ \pm\ -\ \mp;&\text{von }{}^m\beta,\quad \mp\ +\ \pm\ -\\[6pt]
\text{von }{}^m\gamma,\quad \mp\ -\ \pm\ +;&\text{von }{}^m\delta,\quad +\ \mp\ -\ \pm
\end{array}
\]
wobei die oberen Vorzeichen für positives $a$, die unteren für negatives $a$ gelten. Besonders ist folgende Eigenschaft zu beachten: Bezeichnet $m$ einen beliebigen positiven Index, so haben $\alpha^m$ und $\gamma^m$ dasselbe Vorzeichen, wenn $a$ positiv ist. Entgegengesetzte, wenn $a$ negativ ist, und ebenso $\beta^m$ und $\delta^m$; dagegen haben ${}^m\alpha$ und ${}^m\gamma$ oder ${}^m\beta$ und ${}^m\delta$ dasselbe Vorzeichen, wenn $a$ negativ ist, und entgegengesetzte, wenn $a$ positiv ist.

4) In der Bezeichnung von Artikel~27 lässt sich die Größe von $\alpha^m$ usw. bequem so ausdrücken. Man setze
\[
\begin{gathered}
\mp h'=k',\quad \pm h''=k'',\quad \mp h'''=k'''\text{ usw.}\\
\pm h=k,\quad \mp{}'h={}'k,\quad \pm{}''h={}''k\text{ usw.}
\end{gathered}
\]
sodass alle $k',k''$ usw., $k,{}'k$ usw. positive Zahlen sind:
{\small
\setlength{\arraycolsep}{3pt}
\[
\begin{array}{ll}
\alpha^m=\pm[k'',k''',k''''\mathinner{\ldotp\ldotp\ldotp\ldotp} k^{m-1}];&\beta^m=\pm[k'',k''',k''''\mathinner{\ldotp\ldotp\ldotp} {k^m]}\\[6pt]
\gamma^m=\pm[k',k'',k'''\mathinner{\ldotp\ldotp\ldotp\ldotp\ldotp} k^{m-1}];&\delta^m=\pm[k',k'',k'''\mathinner{\ldotp\ldotp\ldotp\ldotp} k^m]\\[6pt]
{}^m\alpha=\pm[k,{}'k,{}''k\mathinner{\ldotp\ldotp\ldotp\ldotp\ldotp\ldotp}{}^{m-1}k];&{}^m\beta=\pm[k,{}'k,{}''k\mathinner{\ldotp\ldotp\ldotp\ldotp}{}^{m-2}k]\\[6pt]
{}^m\gamma=\pm[{}'k,{}''k\mathinner{\ldotp\ldotp\ldotp\ldotp\ldotp\ldotp}{}^{m-1}k];&{}^m\delta=\pm[{}'k,{}''k\mathinner{\ldotp\ldotp\ldotp\ldotp}{}^{m-2}k]
\end{array}
\]
}
Die \textit{Vorzeichen} müssen aber nach den gerade angegebenen Regeln bestimmt werden. Mit diesen Formeln, deren Beweis wir wegen seiner Einfachheit weglassen, lässt sich die Rechnung stets sehr schnell durchführen.
\SourceLeaf{181}

\Needspace{8\baselineskip}
\Article{190}

\textsc{Hilfssatz.} \textit{Bezeichnen $m,\mu,m',n,\nu,n'$ beliebige ganze Zahlen, von denen aber keine der letzten drei $=0$ ist, so behaupte ich: Liegt $\frac{\mu}{\nu}$ streng zwischen $\frac{m}{n}$ und $\frac{m'}{n'}$ und gilt $mn'-nm'={\mp1}$, so ist der Nenner $\nu$ größer als $n$ und $n'$.}

\textit{Beweis.} Offenbar liegt $\mu nn'$ zwischen $\nu mn'$ und $\nu nm'$ und unterscheidet sich daher von jeder Grenze um weniger als die Grenzen voneinander; \textit{das heißt}, es gilt $\nu mn'-\nu nm'>\mu nn'-\nu mn'$ und $>\mu nn'-\nu nm'$, oder $\nu>n'(\mu n-\nu m)$ und $>n(\mu n'-\nu m')$. Da $\mu n-\nu m$ sicher nicht $=0$ ist (sonst wäre $\frac{\mu}{\nu}=\frac{m}{n}$, im Widerspruch zur Voraussetzung), auch nicht $\mu n'-\nu m'=0$ (aus ähnlichem Grund), sondern beide mindestens $=1$ sind, folgt $\nu>n'$ und $>n$. \textit{W.\,z.\,b.\,w.}

Also kann $\nu$ offenbar nicht $=1$ sein; \textit{das heißt}, gilt $mn'-nm'=\pm1$, so kann zwischen den Brüchen $\frac{m}{n}$, $\frac{m'}{n'}$ keine ganze Zahl liegen. Daher kann auch die Null nicht zwischen ihnen liegen, \textit{das heißt}, diese Brüche können keine entgegengesetzten Vorzeichen haben.
